\section{Data and measurement}
\label{sec:data}

This section explains how each submission becomes a data point: where the data came from, which submissions count, how a model is assigned to a company, how the parent is found, and how success is measured.
The paper calls each submission a row.

\subsection{Snapshot}
\label{sec:snapshot}

The snapshot was fetched once from the public API without credentials, by \file{snapshot/fetch.py}.
For every benchmark the API lists, the script stores the benchmark object and its full submission list as gzipped JSON with fixed timestamps.
For the study benchmarks it also stores the attribution record of every submission, cancelled rows included.
An attribution call that returns HTTP 404 is stored as unavailable, and the row's identifier is listed in the freeze record.
After writing the files, the script writes \file{FREEZE.json}, with the fetch time and row counts, and a \file{SHA256SUMS} file covering every data file and \file{FREEZE.json} itself.
Once a freeze exists, the script refuses partial refetches.

The fetch finished at 2026-10-18T13:35:31Z.
It holds 48 benchmarks and 47,525 submissions, of which 32,574 belong to the ten benchmarks in Table~\ref{tab:benchmarks} (\file{snapshot/FREEZE.md}).
Of the attribution calls for those ten benchmarks, 2,150 returned HTTP 404.
At the end of the freeze, a guard checked that every study benchmark holds exactly one attribution record per row, with no duplicates and none missing (\file{docs/JOURNAL.md}).

The loader, \file{snapshot/load.py}, reads only these files and checks them on every open (Appendix~\ref{app:snapshot}).
A submission is flagged as promoted when its status is \field{accepted} and its promotion status is \field{promoted}.

\subsection{Which submissions count}
\label{sec:denominator}

Only submissions that the evaluator actually scored count, since a submission that never ran cannot succeed or fail on its merits.
The paper calls these analysis rows, and \file{outcomes/denominator.py} decides which they are.
The function drops cancelled rows and rows whose status is not accepted, rejected or failed.
It drops accepted rows whose promotion status is set but is not one of promoted, failed or superseded, because their promotion was not final.
It drops rows rejected before evaluation, matched by exact reason strings and patterns in \file{outcomes/denominator.yaml}.
These are rejections on the solver's own claimed score, submissions with no code change, and packaging errors such as an archive path outside the editable paths or an over-long pull-request body.

Failed rows are classed by reason into three groups.
Evaluated failures, where the scoring command ran or the submission's own code failed a gate or a build, are kept.
Infrastructure failures, such as runner setup, artifact transfer, dispatch errors, platform rate limits and interrupted validation jobs, are dropped.
Ambiguous failures, whose text does not say whether the submission caused them, are kept, and sensitivity analysis 9 drops them.
A failed row with no reason counts as ambiguous, and a reason that matches no class stops the pipeline with an error.
The class list was built from the distinct failure reasons in the snapshot, without splitting any outcome by treatment.

\subsection{Which company's model}
\label{sec:labels}

Each submission's note starts with a line naming the model, and the study groups models by the company that makes them, called the model's family.

The declared model is the text after \field{Model:} on the first non-empty line of the note; a note whose first line is anything else has no declared model.
A checked-in table, \file{families/table.yaml}, maps each exact declared string to a family, which is the model's vendor.
The table holds 767 strings.
Of these, 254 map to Unknown and 35 to Multi; the rest map to one of 16 vendor families (\file{families/table.yaml}).

The table was proposed by vendor keyword rules in \file{families/build_table.py} and then corrected by hand where note text gave evidence.
A string naming two or more vendors is Multi, for example ``Claude Opus 5.5 + GPT-6.1 Sol''.
A string naming two models from one vendor maps to that vendor.
A model that a company names as its own maps to that company, as \field{Other:Cognition} for SWE-2 or \field{Other:Cursor} for Composer.
A string maps to Unknown when it names an agent product with no model (``Devin'', ``Codex''), a name whose vendor is ambiguous (``Astra''), or something that is not a model at all.
Strings that name the model being optimised rather than the model doing the work, such as target model weights on the inference benchmarks, are also Unknown.
A row with no \field{Model:} line is Unknown.
Lookup is exact: a declared string absent from the table raises an error, so every analysed string has a recorded decision.

Of the 32,574 study-benchmark rows, 32,386 carry a \field{Model:} line, and 1,453 of those, or 4.5\%, map to Unknown.
Among analysis rows, counting rows with no \field{Model:} line, 5.1\% are Unknown.
Unknown and Multi rows are excluded from the confirmatory tests.
Sensitivity analysis 1 restores Unknown as a family of its own, and the exploratory analysis E1 studies Multi rows.
Within a benchmark, a family with fewer than 30 analysis rows is folded into a single ``Other'' level for the family covariates of the model in Section~\ref{sec:model}.
The \field{Harness:} line, when present, is kept for sensitivity analysis 6, because harness and family are close to collinear.

\paragraph{Label audit.}
A declared label is what the solver typed, so I measured how often the note body agrees with it.
A seeded sample of 200 notes from benchmarks with at least 500 rows was drawn by \file{families/audit/sample.py} (seed 17).
Five fresh-context Claude Opus~5.5 auditors each scored 40 notes as agree, contradict or silent against the \field{Model:} line, ignoring the header itself and verbatim copies of it.
The result was 69 agree, 5 contradict and 126 silent (\file{families/audit/RESULT.md}).
Agreement among non-silent notes was 69 of 74, or 93.2\%, below the pre-registered bar of 95\%.
Four of the five contradictions change the family, three of them because a single-vendor label covered work that the body attributes to more than one vendor.
At the family level, agreement was 70 of 74, or 94.6\%.
The auditors differed on one rule: three of the five batches counted an agent line naming the declared vendor's harness as agreement, and one counted it as silent.
Recoding those 13 notes as silent gives 56 of 61, or 91.8\%, at the model level.
The audit is recorded, not a gate, so the shortfall is a limitation of the labels (Section~\ref{sec:threats}).

\subsection{What each submission built on}
\label{sec:parents}

A submission's parent is the earlier submission it built on.
Participants rarely say so in a machine-readable way, so the study uses three sources, from most to least direct.
The lineage code in \file{lineage/} finds them from three sources, in a fixed order of precedence, and keeps the tree-distance parent for every analysis row as well.

\paragraph{Declared parents.}
When a row's attribution record has \field{declaration} equal to \field{provided} and lists dependencies, those are its declared parents.
Declared identifiers are de-duplicated in order, and identifiers that are unknown or not created before the child are dropped and recorded.
Declared parents are checked against every snapshot row of the benchmark, so a declared parent can be a row that is not itself an analysis row.
One remaining declared parent gives a single-parent row.
Two or more give a multi-source row, which is excluded from the confirmatory tests and studied in E1.

\paragraph{Bases named in the note.}
Many notes name their starting point on a \field{Base:} line.
The code reads that line and resolves it to an earlier submission of the benchmark by identifier prefix, submission commit or promoted commit (\file{lineage/declared.py}).
A resolved base is the parent of a row that has no declared attribution.

\paragraph{Tree-distance parents.}
Each analysis row has a tree: the commit in \field{promotedSourceRef} for promoted rows, and the commit in \field{submissionCommitSha} otherwise.
The build fetches each commit by hash from the benchmark's public repository and records unfetchable ones in \file{missing.json} (Appendix~\ref{app:trees}).
The distance between two trees is the number of added plus removed text lines reported by \texttt{git diff -{}-numstat -{}-minimal -{}-no-renames} over the benchmark's editable paths, with each changed binary file counted as one line.
The candidates are the 50 most recent promoted rows whose promotion finished before the child was created (Appendix~\ref{app:trees}). A solver's own unpromoted rows are not candidates.
The tree parent is the candidate at the smallest distance, with ties going to the most recent row and then to the identifier.

\paragraph{Novel rows.}
A row is novel, and has no tree parent, when the text lines added relative to its nearest candidate exceed half of the text lines in its own tree.
Binary files do not count toward size.
Novel rows are excluded from H1 and H2 and reported descriptively.

\paragraph{Primary parent and freeze.}
The primary parent is the declared one when it exists, then the base named in the note when it resolves to an earlier row, and the tree parent otherwise.
Sensitivity analysis 2 uses the tree parent for every row, which estimates a different quantity: the immediate code parent rather than the declared origin.
The build writes, per benchmark, the tree manifests (path to blob hash), the distance table, the parent links and \file{missing.json} under \file{lineage/frozen/}, with their own \file{SHA256SUMS}.
The analysis rebuilds from these files offline.
The build refuses to overwrite a frozen benchmark when commits it already holds can no longer be fetched.
The final build recorded no unfetchable commit on any study benchmark.

\paragraph{Validation.}
Acceptance test AT4 checks the inferred treatment cell, not the parent identifier, against two labelled sets: declared attribution edges, and notes that name a base in text.
The bar is 85\% agreement on text-declared bases; declared edges are reported but not gated, because only a few dozen of them can be checked.
Agreement was \ATFourTextAgree{} against text-declared bases and \ATFourDeclaredAgree{} against declared edges.
Split by whether the labelled parent was the current best when the child was created, agreement was \ATFourFrontierAgree{} for current-best parents and \ATFourNonFrontierAgree{} for the others.
The parent rule took its final form after a first AT4 run failed before the lock.
That run let a solver's own unpromoted rows compete as candidates even when a different base was named, and agreement on text-declared bases was 1,555 of 3,106, or 50\%, mostly because a solver's own sibling attempt won the distance comparison (\file{docs/JOURNAL.md}).
The revised rule, recorded in the specification before the lock, takes a declared parent first, then a base named in the note, and then the nearest promoted row.

\subsection{The comparison groups}
\label{sec:cells}

A row is cross-solver when its single primary parent was submitted by a different solver.
It is cross-family when the parent's family differs from the child's, with both families known and single-vendor.
Rows with a single parent and two known families fall into four cells: own lineage with the same family, own lineage with a switched family, another solver's work in the same family, and another solver's work in another family.
The primary contrast fixes ``another solver's work'' and varies only the family.
Within a solver, it compares that solver's cross-family attempts with the same solver's same-family attempts, which separates model pollination from people pollination.

\subsection{Measuring success}
\label{sec:outcomes}

The main measure of success is promotion: a submission is promoted when it beats the best score so far by the leaderboard's threshold and becomes the new default starting point.

The outcome code in \file{outcomes/} writes outcome columns only to a sealed directory, and the treatment dataset (\file{analysis/dataset.py}) carries no outcome column.
Outcomes are joined to treatments only inside the confirmatory drivers, after the lock check passes (Section~\ref{sec:lock}).

Rows are ordered by evaluation, using the time promotion finished and falling back to the last update time, because promotions can finish out of creation order.
``Earlier'' in every outcome below means earlier in this order.
The running best before a row is the best promoted score among rows that finished strictly earlier, seeded with the benchmark's baseline score.
\begin{itemize}
  \item Y1, promotion: the row's promotion status is \field{promoted}. Accepted rows whose promotion failed count as not promoted.
  \item Y2, Pareto gain (ECDSA.Fail only): the row has both a qubit count and a Toffoli count, and no earlier row with both counts dominates or equals its point.
  \item Y3, gain size: for a promoted row, the log ratio of the running best before it to its own score, signed so that larger is better under the benchmark's direction.
  \item Large jump: a promotion in the top decile of Y3 within its benchmark.
  \item Plateau break: a promotion that improves the running best by at least 1\% and comes at least 72 hours after the previous promotion of that size, or after the benchmark's opening when there was none.
\end{itemize}

\paragraph{Recomputation.}
Acceptance test AT5 recomputes Y1 in evaluation order from the baseline-seeded running best, applying each benchmark's \field{minScoreImprovementBips}.
Only accepted or rejected rows can promote or move the running best.
Rows whose promotion status is failed or superseded are named exceptions and are removed from the series.
A benchmark whose recomputation disagrees with the recorded promotion status on more than 1\% of its analysis rows is reported as ``promotion rule not reproduced'', and Y1 still comes from the recorded status.
The recomputation reproduced the recorded promotions on \ATFiveReproduced{} benchmarks; the rule was not reproduced on \ATFiveNotReproducedList{}.
On \bench{eigenlabs/flock-challenge-multi/x86}, every row that disagreed in a probe on 17 October predates 25 August 2026, the date of the first rejection citing a 100-basis-point threshold, so the threshold appears to have started part-way through the benchmark.

\subsection{The traceable period}
\label{sec:window}

The chance of finding a submission's code must not depend on whether it succeeded, or the comparison would be skewed.
On ECDSA.Fail, accepted rows lack a submission commit more often than rejected rows, mostly in June, while every promoted row has a promoted commit.
Without a correction, promoted rows would have parents more often than other rows.
For each benchmark, \file{lineage/window.py} groups analysis rows by calendar week, starting on Mondays in UTC.
Walking back from the freeze, it finds the earliest week from which, through the freeze, the share of rows with a recoverable tree differs by at most 5 percentage points between promoted and other rows.
The window starts at that week, and rows before it are reported descriptively only.
Acceptance test AT15 checks the balance inside each window.
On ECDSA.Fail the window starts in the week of \WindowStartEcdsa{}.
